\section{Penutup}
\subsection{Kesimpulan}
Dengan tugas ini, kami berhasil merancang dan mengimplementasikan algoritma block cipher baru bernama OATSIDE. Algoritma ini menggabungkan jaringan Feistel dan lapisan substitusi-permutasi dalam satu putaran, serta menambahkan mekanisme kompresi dan difusi global untuk meningkatkan keamanan. Implementasi OATSIDE dilakukan dalam bahasa Python 3.10, dengan beberapa bagian dioptimalkan menggunakan bahasa C untuk meningkatkan performa.

Berdasarkan hasil pengujian, OATSIDE menunjukkan performa yang baik dalam hal \textit{avalanche effect}, entropi, dan distribusi histogram pada berbagai mode operasi enkripsi. Namun, kami menyadari bahwa masih ada ruang untuk perbaikan dan optimasi lebih lanjut, terutama dalam hal performa dan keamanan algoritma.

\subsection{Saran}
Sebaiknya pengembangan selanjutnya pakai bahasa yang lebih sigma seperti C atau Rust, jangan pakai Python karena lemot. Selain itu, kami punya hipotesis bahwa ada beberapa hal dari skema \textit{block cipher}-nya yang \textit{by design} kurang optimal secara performa maupun keamanan, jadi sebaiknya dicoba untuk diubah-ubah dan diuji lagi. 

\subsection{Lampiran}
Tautan dokumentasi API: \href{https://brianhadianstei23.github.io/13523048-13523094-135230115-tubes2-IF4020/}{https://brianhadianstei23.github.io/13523048-13523094-135230115-tubes2-IF4020/}
\\ 
Video demo: 
\pagebreak